% ---
% titre: Multiplications et divisions lacunaires (séries A, B et C)
% theme: NO
% chapitre: Nombres naturels et décimaux
% sous_chapitre: Calcul réfléchi et mental
% niveau: [9e]
% type: EX
% tags: [c-livrets, c-multiplication, c-division, f-individuel,f-maison,f-station,p-entrainement,p-remediation]
% difficulte: 2
% source: latex
% ---

\header{Multiplications et divisions lacunaires (série A)}

Trouve le nombre manquant dans chaque égalité pour compléter les séries ci-dessous.

\vfill
\def\x{0.7}
\def\arraystretch{2}
\begin{tabularx}{1\linewidth}{X r c l c l @{\hspace{15pt}}|@{\hspace{15pt}} X r c l c l @{\hspace{15pt}}|@{\hspace{15pt}} X r c l c l}
1) &\sol{\bl{\x}}{6} & $\times$ & 4 & = & 24 & 21) &\sol{\bl{\x}}{132} & $\div$ & 12 & = & 11 & 41) &\sol{\bl{\x}}{7} & $\times$ & 6 & = & 42 \\
2) & \sol{\bl{\x}}{3} & $\times$ & 5 & = & 15 & 22) & \sol{\bl{\x}}{22} & $\div$ & 11 & = & 2 & 42) & 48 & $\div$ & \sol{\bl{\x}}{8} & = & 6 \\
3) & \sol{\bl{\x}}{10} & $\times$ & 8 & = & 80 & 23) & 27 & $\div$ & 9 & = & \sol{\bl{\x}}{3} & 43) & 9 & $\times$ & \sol{\bl{\x}}{7} & = & 63 \\
4) & \sol{\bl{\x}}{9} & $\times$ & 3 & = & 27 & 24) & 32 & $\div$ & 8 & = & \sol{\bl{\x}}{4} & 44) & \sol{\bl{\x}}{72} & $\div$ & 8 & = & 9 \\
5) & \sol{\bl{\x}}{6} & $\times$ & 3 & = & 18 & 25) & \sol{\bl{\x}}{88} & $\div$ & 11 & = & 8 & 45) & 5 & $\times$ & 7 & = & \sol{\bl{\x}}{35} \\
6) & 4 & $\times$ & 11 & = & \sol{\bl{\x}}{44} & 26) & 24 & $\div$ & \sol{\bl{\x}}{3} & = & 8 & 46) & 121 & $\div$ & 11 & = & \sol{\bl{\x}}{11} \\
7) & \sol{\bl{\x}}{12} & $\times$ & 10 & = & 120 & 27) & 33 & $\div$ & 11 & = & \sol{\bl{\x}}{3} & 47) & \sol{\bl{\x}}{8} & $\times$ & 12 & = & 96 \\
8) & 8 & $\times$ & \sol{\bl{\x}}{9} & = & 72 & 28) & 63 & $\div$ & \sol{\bl{\x}}{9} & = & 7 & 48) & 63 & $\div$ & \sol{\bl{\x}}{7} & = & 9 \\
9) & 5 & $\times$ & \sol{\bl{\x}}{9} & = & 45 & 29) & 56 & $\div$ & 7 & = & \sol{\bl{\x}}{8} & 49) & 8 & $\times$ & 9 & = & \sol{\bl{\x}}{72} \\
10) & 11 & $\times$ & 11 & = & \sol{\bl{\x}}{121} & 30) & 77 & $\div$ & \sol{\bl{\x}}{7} & = & 11 & 50) & \sol{\bl{\x}}{44} & $\div$ & 4 & = & 11 \\
11) & \sol{\bl{\x}}{5} & $\times$ & 10 & = & 50 & 31) & 27 & $\div$ & 9 & = & \sol{\bl{\x}}{3} & 51) & 6 & $\times$ & \sol{\bl{\x}}{9} & = & 54 \\
12) & 12 & $\times$ & 9 & = & \sol{\bl{\x}}{108} & 32) & \sol{\bl{\x}}{36} & $\div$ & 4 & = & 9 & 52) & 100 & $\div$ & \sol{\bl{\x}}{10} & = & 10 \\
13) & 2 & $\times$ & \sol{\bl{\x}}{8} & = & 16 & 33) & \sol{\bl{\x}}{8} & $\div$ & 2 & = & 4 & 53) & \sol{\bl{\x}}{11} & $\times$ & 3 & = & 33 \\
14) & 12 & $\times$ & 12 & = & \sol{\bl{\x}}{144} & 34) & 16 & $\div$ & \sol{\bl{\x}}{8} & = & 2 & 54) & 45 & $\div$ & 5 & = & \sol{\bl{\x}}{9} \\
15) & 11 & $\times$ & 12 & = & \sol{\bl{\x}}{132} & 35) & \sol{\bl{\x}}{27} & $\div$ & 3 & = & 9 & 55) & 7 & $\times$ & \sol{\bl{\x}}{12} & = & 84 \\
16) & 11 & $\times$ & 8 & = & \sol{\bl{\x}}{88} & 36) & 20 & $\div$ & \sol{\bl{\x}}{4} & = & 5 & 56) & \sol{\bl{\x}}{72} & $\div$ & 6 & = & 12 \\
17) & 10 & $\times$ & 8 & = & \sol{\bl{\x}}{80} & 37) & \sol{\bl{\x}}{45} & $\div$ & 9 & = & 5 & 57) & 9 & $\times$ & 6 & = & \sol{\bl{\x}}{54} \\
18) & 10 & $\times$ & \sol{\bl{\x}}{9} & = & 90 & 38) & 96 & $\div$ & 8 & = & \sol{\bl{\x}}{12} & 58) & 132 & $\div$ & 12 & = & \sol{\bl{\x}}{11} \\
19) & 12 & $\times$ & 9 & = & \sol{\bl{\x}}{108} & 39) & 33 & $\div$ & 11 & = & \sol{\bl{\x}}{3} & 59) & \sol{\bl{\x}}{11} & $\times$ & 8 & = & 88 \\
20) & 10 & $\times$ & \sol{\bl{\x}}{3} & = & 30 & 40) & 12 & $\div$ & \sol{\bl{\x}}{4} & = & 3 & 60) & 24 & $\div$ & \sol{\bl{\x}}{3} & = & 8 \\
\end{tabularx}
\def\arraystretch{1}

\newpage
\header{Multiplications et divisions lacunaires (série B)}

Trouve le nombre manquant dans chaque égalité pour compléter les séries ci-dessous.

\vfill
\def\x{0.7}
\def\arraystretch{2}
\begin{tabularx}{1\linewidth}{X r c l c l @{\hspace{15pt}}|@{\hspace{15pt}} X r c l c l @{\hspace{15pt}}|@{\hspace{15pt}} X r c l c l}
1) & 6 & $\times$ & \sol{\bl{\x}}{5} & = & 30 & 21) & 70 & $\div$ & \sol{\bl{\x}}{7} & = & 10 & 41) & \sol{\bl{\x}}{8} & $\times$ & 9 & = & 72 \\
2) & 7 & $\times$ & \sol{\bl{\x}}{11} & = & 77 & 22) & \sol{\bl{\x}}{40} & $\div$ & 10 & = & 4 & 42) & 56 & $\div$ & \sol{\bl{\x}}{7} & = & 8 \\
3) & 2 & $\times$ & 5 & = & \sol{\bl{\x}}{10} & 23) & 72 & $\div$ & \sol{\bl{\x}}{6} & = & 12 & 43) & 4 & $\times$ & \sol{\bl{\x}}{12} & = & 48 \\
4) & \sol{\bl{\x}}{11} & $\times$ & 2 & = & 22 & 24) & 24 & $\div$ & \sol{\bl{\x}}{8} & = & 3 & 44) & \sol{\bl{\x}}{66} & $\div$ & 11 & = & 6 \\
5) & \sol{\bl{\x}}{4} & $\times$ & 9 & = & 36 & 25) & 108 & $\div$ & 9 & = & \sol{\bl{\x}}{12} & 45) & 8 & $\times$ & 8 & = & \sol{\bl{\x}}{64} \\
6) & 2 & $\times$ & \sol{\bl{\x}}{4} & = & 8 & 26) & 48 & $\div$ & \sol{\bl{\x}}{6} & = & 8 & 46) & 90 & $\div$ & 9 & = & \sol{\bl{\x}}{10} \\
7) & 10 & $\times$ & \sol{\bl{\x}}{11} & = & 110 & 27) & \sol{\bl{\x}}{14} & $\div$ & 7 & = & 2 & 47) & \sol{\bl{\x}}{9} & $\times$ & 7 & = & 63 \\
8) & 5 & $\times$ & 8 & = & \sol{\bl{\x}}{40} & 28) & \sol{\bl{\x}}{132} & $\div$ & 12 & = & 11 & 48) & 108 & $\div$ & \sol{\bl{\x}}{9} & = & 12 \\
9) & 7 & $\times$ & \sol{\bl{\x}}{4} & = & 28 & 29) & \sol{\bl{\x}}{33} & $\div$ & 3 & = & 11 & 49) & 5 & $\times$ & 9 & = & \sol{\bl{\x}}{45} \\
10) & \sol{\bl{\x}}{4} & $\times$ & 7 & = & 28 & 30) & 110 & $\div$ & 10 & = & \sol{\bl{\x}}{11} & 50) & \sol{\bl{\x}}{36} & $\div$ & 3 & = & 12 \\
11) & 9 & $\times$ & \sol{\bl{\x}}{2} & = & 18 & 31) & \sol{\bl{\x}}{35} & $\div$ & 7 & = & 5 & 51) & 11 & $\times$ & \sol{\bl{\x}}{9} & = & 99 \\
12) & 9 & $\times$ & \sol{\bl{\x}}{6} & = & 54 & 32) & 72 & $\div$ & \sol{\bl{\x}}{12} & = & 6 & 52) & 84 & $\div$ & \sol{\bl{\x}}{12} & = & 7 \\
13) & \sol{\bl{\x}}{2} & $\times$ & 2 & = & 4 & 33) & 90 & $\div$ & 9 & = & \sol{\bl{\x}}{10} & 53) & \sol{\bl{\x}}{11} & $\times$ & 6 & = & 66 \\
14) & 12 & $\times$ & \sol{\bl{\x}}{5} & = & 60 & 34) & \sol{\bl{\x}}{132} & $\div$ & 11 & = & 12 & 54) & 96 & $\div$ & 8 & = & \sol{\bl{\x}}{12} \\
15) & 11 & $\times$ & \sol{\bl{\x}}{8} & = & 88 & 35) & \sol{\bl{\x}}{90} & $\div$ & 10 & = & 9 & 55) & 3 & $\times$ & \sol{\bl{\x}}{12} & = & 36 \\
16) & 7 & $\times$ & 10 & = & \sol{\bl{\x}}{70} & 36) & 30 & $\div$ & \sol{\bl{\x}}{6} & = & 5 & 56) & \sol{\bl{\x}}{72} & $\div$ & 9 & = & 8 \\
17) & 7 & $\times$ & \sol{\bl{\x}}{3} & = & 21 & 37) & 4 & $\div$ & \sol{\bl{\x}}{2} & = & 2 & 57) & 7 & $\times$ & 12 & = & \sol{\bl{\x}}{84} \\
18) & 6 & $\times$ & \sol{\bl{\x}}{6} & = & 36 & 38) & 30 & $\div$ & \sol{\bl{\x}}{5} & = & 6 & 58) & 110 & $\div$ & 11 & = & \sol{\bl{\x}}{10} \\
19) & \sol{\bl{\x}}{3} & $\times$ & 6 & = & 18 & 39) & \sol{\bl{\x}}{77} & $\div$ & 7 & = & 11 & 59) & \sol{\bl{\x}}{10} & $\times$ & 4 & = & 40 \\
20) & \sol{\bl{\x}}{6} & $\times$ & 5 & = & 30 & 40) & \sol{\bl{\x}}{27} & $\div$ & 9 & = & 3 & 60) & 18 & $\div$ & \sol{\bl{\x}}{3} & = & 6 \\
\end{tabularx}
\def\arraystretch{1}

\newpage
\header{Multiplications et divisions lacunaires (série C)}

Trouve le nombre manquant dans chaque égalité pour compléter les séries ci-dessous.

\vfill
\def\x{0.7}
\def\arraystretch{2}
\begin{tabularx}{1\linewidth}{X r c l c l @{\hspace{15pt}}|@{\hspace{15pt}} X r c l c l @{\hspace{15pt}}|@{\hspace{15pt}} X r c l c l}
1) & \sol{\bl{\x}}{7} & $\times$ & 12 & = & 84 & 21) & \sol{\bl{\x}}{36} & $\div$ & 6 & = & 6 & 41) & \sol{\bl{\x}}{7} & $\times$ & 11 & = & 77 \\
2) & \sol{\bl{\x}}{3} & $\times$ & 8 & = & 24 & 22) & 55 & $\div$ & 11 & = & \sol{\bl{\x}}{5} & 42) & 42 & $\div$ & \sol{\bl{\x}}{7} & = & 6 \\
3) & 10 & $\times$ & \sol{\bl{\x}}{12} & = & 120 & 23) & 80 & $\div$ & \sol{\bl{\x}}{10} & = & 8 & 43) & 6 & $\times$ & \sol{\bl{\x}}{12} & = & 72 \\
4) & 8 & $\times$ & 2 & = & \sol{\bl{\x}}{16} & 24) & \sol{\bl{\x}}{18} & $\div$ & 3 & = & 6 & 44) & \sol{\bl{\x}}{96} & $\div$ & 12 & = & 8 \\
5) & \sol{\bl{\x}}{7} & $\times$ & 12 & = & 84 & 25) & 18 & $\div$ & \sol{\bl{\x}}{2} & = & 9 & 45) & 9 & $\times$ & 8 & = & \sol{\bl{\x}}{72} \\
6) & 8 & $\times$ & 7 & = & \sol{\bl{\x}}{56} & 26) & 90 & $\div$ & \sol{\bl{\x}}{9} & = & 10 & 46) & 63 & $\div$ & 9 & = & \sol{\bl{\x}}{7} \\
7) & 10 & $\times$ & 4 & = & \sol{\bl{\x}}{40} & 27) & \sol{\bl{\x}}{36} & $\div$ & 4 & = & 9 & 47) & \sol{\bl{\x}}{12} & $\times$ & 5 & = & 60 \\
8) & 9 & $\times$ & 7 & = & \sol{\bl{\x}}{63} & 28) & 15 & $\div$ & \sol{\bl{\x}}{5} & = & 3 & 48) & 88 & $\div$ & \sol{\bl{\x}}{8} & = & 11 \\
9) & \sol{\bl{\x}}{7} & $\times$ & 7 & = & 49 & 29) & \sol{\bl{\x}}{77} & $\div$ & 11 & = & 7 & 49) & 4 & $\times$ & 7 & = & \sol{\bl{\x}}{28} \\
10) & 3 & $\times$ & \sol{\bl{\x}}{9} & = & 27 & 30) & \sol{\bl{\x}}{14} & $\div$ & 7 & = & 2 & 50) & \sol{\bl{\x}}{90} & $\div$ & 10 & = & 9 \\
11) & \sol{\bl{\x}}{11} & $\times$ & 11 & = & 121 & 31) & 32 & $\div$ & \sol{\bl{\x}}{8} & = & 4 & 51) & 12 & $\times$ & \sol{\bl{\x}}{11} & = & 132 \\
12) & 11 & $\times$ & \sol{\bl{\x}}{7} & = & 77 & 32) & 60 & $\div$ & 12 & = & \sol{\bl{\x}}{5} & 52) & 77 & $\div$ & \sol{\bl{\x}}{11} & = & 7 \\
13) & 5 & $\times$ & \sol{\bl{\x}}{11} & = & 55 & 33) & \sol{\bl{\x}}{88} & $\div$ & 8 & = & 11 & 53) & \sol{\bl{\x}}{12} & $\times$ & 9 & = & 108 \\
14) & 11 & $\times$ & \sol{\bl{\x}}{10} & = & 110 & 34) & \sol{\bl{\x}}{49} & $\div$ & 7 & = & 7 & 54) & 40 & $\div$ & 8 & = & \sol{\bl{\x}}{5} \\
15) & \sol{\bl{\x}}{9} & $\times$ & 8 & = & 72 & 35) & 12 & $\div$ & 4 & = & \sol{\bl{\x}}{3} & 55) & 6 & $\times$ & \sol{\bl{\x}}{8} & = & 48 \\
16) & \sol{\bl{\x}}{12} & $\times$ & 5 & = & 60 & 36) & \sol{\bl{\x}}{24} & $\div$ & 12 & = & 2 & 56) & \sol{\bl{\x}}{77} & $\div$ & 7 & = & 11 \\
17) & 10 & $\times$ & 3 & = & \sol{\bl{\x}}{30} & 37) & \sol{\bl{\x}}{60} & $\div$ & 5 & = & 12 & 57) & 11 & $\times$ & 3 & = & \sol{\bl{\x}}{33} \\
18) & 6 & $\times$ & \sol{\bl{\x}}{10} & = & 60 & 38) & \sol{\bl{\x}}{54} & $\div$ & 9 & = & 6 & 58) & 60 & $\div$ & 12 & = & \sol{\bl{\x}}{5} \\
19) & \sol{\bl{\x}}{2} & $\times$ & 3 & = & 6 & 39) & 60 & $\div$ & 5 & = & \sol{\bl{\x}}{12} & 59) & \sol{\bl{\x}}{12} & $\times$ & 2 & = & 24 \\
20) & 7 & $\times$ & 7 & = & \sol{\bl{\x}}{49} & 40) & 96 & $\div$ & 8 & = & \sol{\bl{\x}}{12} & 60) & 36 & $\div$ & \sol{\bl{\x}}{9} & = & 4 \\
\end{tabularx}
\def\arraystretch{1}
